% ==========================================================================
%  preamble.tex --- theme, macros and theorem environments for the talk.
%
%  Deliberately sober: no navigation bar, no shadows, no gradients. The
%  palette is a single slate blue plus a muted accent, so that displayed
%  mathematics rather than decoration carries the eye.
% ==========================================================================

\usetheme{default}
\usecolortheme{default}
\setbeamertemplate{navigation symbols}{}
\setbeamertemplate{itemize items}[circle]
\setbeamertemplate{enumerate items}[default]
\setbeamertemplate{sections/subsections in toc}[circle]

% ---------- fonts ---------------------------------------------------------
% Palatino body and matching math, as in the paper itself.
\usepackage[T1]{fontenc}
\usepackage[utf8]{inputenc}
\usepackage{newpxtext}
\usepackage{amsmath,amssymb,mathtools}
\usepackage{newpxmath}
\usepackage{microtype}
\usefonttheme{professionalfonts}

\usepackage{booktabs}
\usepackage{array}
\usepackage{tikz}
\usetikzlibrary{arrows.meta,positioning,calc,fit,backgrounds}

% ---------- palette -------------------------------------------------------
\definecolor{Slate}{RGB}{38,62,94}
\definecolor{SlateLight}{RGB}{232,238,246}
\definecolor{Accent}{RGB}{150,40,40}
\definecolor{Ink}{RGB}{30,30,32}
\definecolor{Muted}{RGB}{110,116,124}

\setbeamercolor{normal text}{fg=Ink}
\setbeamercolor{structure}{fg=Slate}
\setbeamercolor{frametitle}{fg=Slate,bg=}
\setbeamercolor{title}{fg=Slate}
\setbeamercolor{section in toc}{fg=Ink}
\setbeamercolor{alerted text}{fg=Accent}

\setbeamerfont{frametitle}{size=\large,series=\bfseries}
\setbeamerfont{title}{size=\LARGE,series=\bfseries}
\setbeamerfont{author}{size=\normalsize}

% Frame title with a thin rule underneath -- enough structure to orient the
% eye, not enough to compete with the mathematics.
\setbeamertemplate{frametitle}{%
  \vspace{0.55em}%
  \begin{minipage}{\textwidth}
    \raggedright
    {\usebeamerfont{frametitle}\usebeamercolor[fg]{frametitle}\insertframetitle}\par
    \vspace{0.22em}%
    {\color{Slate}\rule{\textwidth}{0.6pt}}\par
  \end{minipage}%
  \vspace{0.15em}%
}

% Footline: section on the left, frame number on the right.
\setbeamertemplate{footline}{%
  \hbox{\begin{beamercolorbox}[wd=\paperwidth,ht=2.4ex,dp=1.2ex,leftskip=1.2em,rightskip=1.2em]{}%
    {\scriptsize\color{Muted}\insertsection\hfill\insertframenumber/\inserttotalframenumber}%
  \end{beamercolorbox}}%
}

% ---------- theorem environments -----------------------------------------
\setbeamertemplate{blocks}[rounded][shadow=false]
\setbeamercolor{block title}{bg=SlateLight,fg=Slate}
\setbeamercolor{block body}{bg=SlateLight!40,fg=Ink}
\setbeamercolor{block title alerted}{bg=Accent!12,fg=Accent}
\setbeamercolor{block body alerted}{bg=Accent!5,fg=Ink}

\setbeamertemplate{theorems}[numbered]
\theoremstyle{definition}
\newtheorem{defn}{Definition}
\newtheorem{prop}[theorem]{Proposition}
\newtheorem{obs}[theorem]{Observation}

% A quieter "Proof." head than beamer's default block.
\newenvironment{pf}[1][Proof]{%
  \par\vspace{0.2em}\noindent{\color{Slate}\scshape #1.}\ \ignorespaces
}{\par\vspace{0.2em}}

% Step markers inside multi-part proofs.
\newcommand{\step}[1]{\medskip\noindent{\color{Slate}\bfseries #1}\quad}

% ---------- project notation (identical to the paper) ---------------------
\newcommand{\agents}{N}
\newcommand{\items}{M}
\newcommand{\alloc}{\mathcal{A}}
\newcommand{\allocB}{\mathcal{B}}
\newcommand{\bundle}[1]{A_{#1}}
\newcommand{\subsidy}{p}
\newcommand{\optsubsidy}{p^{*}}
\newcommand{\cost}{c}
\newcommand{\N}{\mathbb{N}}
\newcommand{\R}{\mathbb{R}}
\newcommand{\Z}{\mathbb{Z}}
\newcommand{\set}[1]{\left\{#1\right\}}
\newcommand{\abs}[1]{\left\lvert#1\right\rvert}

% Emphasis for the quantity under discussion.
\newcommand{\hl}[1]{{\color{Accent}#1}}
\newcommand{\cyc}{\textnormal{(R1)}}

% Reference tags, written out rather than BibTeX'd: a talk needs the names
% visible on the slide, not a numeral to look up.
\newcommand{\refr}[1]{{\footnotesize\color{Muted}[#1]}}
